We evaluate the framework in four parts. We first calibrate every cost term on
real hardware and select the deployed models from the measured accuracy-energy
frontier. We then compare the proposed scheduler against four baselines that
draw from the same feasible action set, reporting the comparison as a
quality-fairness trade-off rather than a single operating point. We validate
the theoretical guarantees against a certified offline optimum and long-horizon
stability runs. Throughout, we report the boundaries of every empirical
advantage explicitly; Section~\ref{EV_LIM} collects them.

\subsection{Hardware Calibration and Model Selection}\label{EV_CALIB}
We profile seven vision-language models on RESISC45 (1300 images, 19 to 41 per
category), each at 4-bit quantization on an NVIDIA Jetson Orin NX in its 15\,W
power mode. Latency is timed per image, power is sampled at 50\,ms intervals
with \texttt{tegrastats}, and energy integrates power over the inference
window. Fig.~\ref{fig:tradeoff} places every model on the accuracy-energy
plane. Four models lie on the Pareto frontier: Qwen2-VL-2B (0.394 accuracy,
3.93\,J), Qwen2.5-VL-3B (0.456, 5.72\,J), Qwen2.5-VL-7B (0.549, 9.30\,J), and
InternVL3-8B, which reaches the highest accuracy (0.562) at the highest energy
(15.00\,J). The other three are dominated: SmolVLM2-2.2B (0.310),
InternVL2.5-4B (0.548), and gemma-3-4b (0.538, at 37.79\,J).

We adopt the Qwen family as the deployment backbone for an architectural
reason, not a per-point accuracy judgement. Qwen2-VL-2B, Qwen2.5-VL-3B, and
Qwen2.5-VL-7B share one architecture and one activation interface, so they form
a single three-tier ladder that maps onto the deployment tiers of this work: a
light single-satellite model, a heavier single-satellite model, and a
split-only model. Their uniform transformer-block structure keeps the
activation dimension constant at any partition point, which is the precondition
for the two-stage split. InternVL3-8B is an isolated 8B model of a different
architecture: it does not compose with the on-board 2B and 3B tiers into one
routable family, its larger size is even less feasible on a single satellite,
and splitting it is heavier; its accuracy lead therefore does not change the
deployment choice. InternVL2.5-4B is within statistical noise of Qwen2.5-VL-7B
(0.548 versus 0.549). SmolVLM2-2.2B is the weakest model and is dominated by
the smaller Qwen2-VL-2B, which shows that a model's frontier position is set by
architecture and training rather than by parameter count alone.

The split pipeline is required by a hardware constraint that does not depend on
any accuracy comparison. A 4-bit Qwen2.5-VL-7B model, with its vision encoder
and key-value cache, does not fit in the memory left on an 8\,GB satellite
after the on-board sensing pipeline, so the two-satellite split is the only way
to run the 7B tier at all.

\begin{table}[!t]
\centering\small
\caption{\label{tab_calib}Measured per-configuration cost on Jetson Orin NX
(RESISC45, $n{=}1300$, \texttt{Q4\_K\_M}, 15\,W mode). The peak bound
$\hat{P}_{im}^{\text{CP}}$ is the maximum observed per-inference peak, a
conservative upper bound for the hard per-slot cap. The 7B tier is split-only
under the 8\,GB satellite budget; $N_{im}$ is the number of $\tau{=}1$\,s
slots.}
\begin{tabular}{lccccc}
\toprule
Config & $Q_{im}$ & $\mathbb{E}[T_{im}]$(s) & $\mathbb{E}[E_{im}]$(J) & $\hat{P}_{im}^{\text{CP}}$(W) & $N_{im}$ \\
\midrule
Qwen2-VL-2B   & 0.394 & 1.24 & 3.93 & 3.75 & 2 \\
Qwen2.5-VL-3B & 0.456 & 1.74 & 5.72 & 4.02 & 2 \\
Qwen2.5-VL-7B & 0.549 & 2.75 & 9.30 & 6.37 & 3 \\
\bottomrule
\end{tabular}
\end{table}

\begin{figure}[!t]
\centering
\includegraphics[width=\columnwidth]{quality_energy_tradeoff_RESISC45.pdf}
\caption{Accuracy-energy landscape of the seven profiled models (RESISC45,
$n{=}1300$, all 4-bit). The Pareto frontier holds four models; the Qwen
2B/3B/7B ladder occupies its low-energy segment and forms the deployed
three-tier family.}
\label{fig:tradeoff}
\end{figure}

\subsection{Peak-Power Calibration}\label{EV_PEAK}
The instantaneous power cap is the second structural constraint. We model it as
a shared power budget $P^{\text{cap}}{=}15$\,W, taking the Jetson 15\,W
envelope as the nominal value; the platform baseline $P_{\text{base}}$ and the
inference peak draw from this budget. This is an abstraction: a real satellite
separates the platform power system from the compute payload, and a
higher-fidelity power model is left to future work. We calibrate the peak bound
$\hat{P}_{im}^{\text{CP}}$ as the maximum per-inference peak observed over the
benchmark, a conservative upper bound appropriate for a hard constraint that
must never be violated. Among the deployed configurations only the 7B tier
shows a sharp transient: its peak reaches 6.37\,W against a 3.77\,W median (a
$1.9\times$ ratio), traced to a sustained compute-bound prefill phase, whereas
the 2B and 3B peaks (3.75 and 4.02\,W) sit close to their own high quantiles.
The 7B peak plus a few-watt platform baseline approaches the 15\,W budget, so a
scheduler that commits the split without checking the peak can drive the
combined draw past the cap. This is the empirical basis for the per-slot peak
shield.

\subsection{Simulation Setup}\label{EV_SETUP}
Each slot is $\tau{=}1$\,s; one orbit is 5700 slots (60\% sunlit, 40\% eclipse)
and we run three orbits unless noted. The constellation has four computing
satellites and six sensing sources with heterogeneous quality profiles taken
from the per-category RESISC45 measurements. The platform baseline continuously
draws $P_{\text{base}}$ from the battery; in eclipse it is borne entirely by
the battery and forms the bulk of the eclipse energy budget. We take
$P_{\text{base}}{=}7$\,W as a nominal value in the reported CubeSat range
(3--10\,W for attitude control and housekeeping) and sweep it in
Section~\ref{EV_SENS}; the solar panel ($P^{\text{solar}}{=}13$\,W) yields an
orbit-average harvest just above the platform draw, leaving a thin margin for
paced inference, and the battery (18\,kJ) is sized to carry the platform
through eclipse. All policies draw from one shared feasible set (the split
included, the ISL handoff check, the per-slot energy check, and the peak
check), so they differ only in the selection rule; quality ties are broken
uniformly at random. Results are means over 10 seeds. The simulation uses a
serial scheduling implementation that commits one decision per slot,
corresponding to a centralized executor in which a lead satellite collects
state and broadcasts the decision. The abstraction is identical for every
policy, so it does not affect the fairness of the comparison, but it means
throughput does not scale linearly with constellation size; the formulation
itself admits concurrent execution, of which the serial implementation is a
conservative special case.

We compare against four baselines. Greedy-Q takes the highest-quality feasible
action, reaching for the 7B split whenever two satellites are free. Greedy-E
takes the cheapest feasible action. Random selects uniformly at random from the
feasible set. Static runs a fixed 3B standalone configuration on the
highest-battery satellite. We treat the offline MILP solely as a clairvoyant
upper bound for the optimality-gap study (Section~\ref{EV_GAP}): it is
non-causal and intractable at constellation scale.

\subsection{Constrained-Fair Comparison}\label{EV_MULTI}
Fig.~\ref{fig:frontier} and Table~\ref{tab_split} report the four-satellite
comparison on the quality-fairness plane. The proposed scheduler is not a
single point but a curve swept by the Lyapunov weight $V$: small $V$ favours
the max-min objective (min-fill up to 0.175), large $V$ favours total
throughput (up to 3850), and we pin no single $V$. The baselines are fixed
points. The reading is threefold. First, the proposed scheduler is the only
policy that is simultaneously deployable and able to run the 7B split: it holds
downtime at 0.4\% across the whole $V$ range while running 179 to 206 splits,
whereas the two other deployable policies (Greedy-E and Static) never split.
Second, at small $V$ it exceeds the fair baseline Greedy-E on min-fill (0.175
versus 0.156, a $5\sigma$ gap) at comparable total quality (3254 versus 3193, a
difference within $1.6\sigma$, not significant). Third, the policies that reach
both high fairness and high throughput (Greedy-Q at min-fill 0.398, Random at
0.354) do so only by blacking out (7.0\% and 5.3\% downtime), which violates
the uptime service-level agreement (SLA); the deployability partition is
invariant for any SLA threshold in the wide band $(0.8\%, 5.3\%)$. No policy
reaches high fairness and high throughput while remaining deployable, which is
the quality-fairness conflict made visible as the empty upper region of
Fig.~\ref{fig:frontier}. Static reaches higher raw throughput (5102) at low
fairness (0.119); the proposed scheduler trades this throughput for the ability
to navigate to the fair region and to run the top tier.

\begin{table}[!t]
\centering\footnotesize
\setlength{\tabcolsep}{3pt}
\caption{\label{tab_split}Four-satellite comparison (three orbits, 10-seed
means; standard deviations below 0.007 for min-fill and 70 for total quality).
All policies draw from the same shared feasible set. The proposed scheduler is
reported at three values of the Lyapunov weight $V$ to show the
fairness-throughput knob; no single $V$ is selected (Fig.~\ref{fig:frontier}).
Downtime above the uptime SLA marks a policy undeployable.}
\begin{tabular}{lcccccc}
\toprule
Policy & \makecell{Peak\\viol.} & \makecell{Down\\time} & \makecell{Serv.\\rate} & \makecell{7B\\splits} & \makecell{Min-\\fill} & \makecell{Total\\qual.} \\
\midrule
Proposed ($V{=}1$)   & 0 & 0.4\% & 0.525 & $206$ & $\mathbf{0.175}$ & $3254$ \\
Proposed ($V{=}10$)  & 0 & 0.4\% & 0.526 & $201$ & $0.169$ & $3311$ \\
Proposed ($V{=}100$) & 0 & 0.4\% & 0.534 & $179$ & $0.124$ & $3636$ \\
\midrule
Greedy-Q & 0 & \textbf{7.0\%} & 0.765 & 5299 & 0.398 & 6083 \\
Greedy-E & 0 & 0.8\% & 0.780 & 0 & 0.156 & 3193 \\
Random   & 0 & \textbf{5.3\%} & 0.778 & 3333 & 0.354 & 5022 \\
Static   & 0 & 0.8\% & 0.780 & 0 & 0.119 & 5102 \\
\bottomrule
\end{tabular}
\end{table}

\begin{figure}[!t]
\centering
\includegraphics[width=\columnwidth]{tradeoff_frontier.pdf}
\caption{The quality-fairness trade-off at four satellites (10-seed means).
The proposed scheduler is a curve swept by $V$, from the fairness end (small
$V$) to the throughput end; no single $V$ is selected. At small $V$ its
min-fill exceeds Greedy-E significantly ($5\sigma$) at total quality within
noise ($1.6\sigma$). Greedy-Q and Random (hollow) reach the upper region only
by blacking out, violating the SLA; the partition holds for any SLA in
$(0.8\%,5.3\%)$. The empty deployable upper-right region is the
quality-fairness conflict.}
\label{fig:frontier}
\end{figure}

\subsection{Capability and Routing}\label{EV_ROUTE}
The heterogeneity of the sources makes model choice a routing problem rather
than a size ranking. To avoid selecting categories after seeing the data, we
fixed a test in advance: a category requires the split-only 7B model if 3B
scores at most 0.35, 7B scores at least 0.70, and the difference is significant
under a two-sided Fisher exact test after Benjamini-Hochberg correction across
all 45 categories. Five categories pass. Four of them, however, are
non-monotone in model size: the smaller 2B model scores 0.79 to 0.89, 3B
collapses to 0.04 to 0.34, and 7B recovers to 0.78 to 0.83. This is a
category-specific weakness of the 3B model, and because 2B also runs on a
single satellite, the best single-satellite option already serves these four.
Under the stricter criterion that both single-satellite models fail
($\max(\text{2B},\text{3B})\le 0.35$) while 7B reaches 0.70, which we mark as a
deviation from the pre-registered test, only one category (mountain) qualifies,
and it passes Benjamini-Hochberg but not Bonferroni. The aggregate accuracy
gain of 7B over 3B is 0.093, and the per-category sample is small (19 to 41
images, $95\%$ interval about $\pm0.18$), so we rest the case for the split on
the memory constraint rather than on any single category. The value of the data
is that model size alone is a poor guide to configuration: most categories are
served by a single-satellite model, several on which 3B fails are recovered by
the smaller 2B, and the scheduler routes on the measured per-source quality
rather than on size.

\subsection{The Fairness--Throughput Conflict}\label{EV_VCONF}
Fig.~\ref{fig:vconf} traces the two objectives against $V$. Total delivered
quality rises with $V$, so the $O(1/V)$ improvement of Theorem~\ref{thm_gap}
holds in the multi-satellite setting. Min-fill, in contrast, falls with $V$.
This reverses its single-satellite behaviour, where min-quality rises with $V$
because there is no fairness conflict; here the heterogeneous sources create
one. The reversal does not violate Theorem~\ref{thm_gap}, which bounds total
quality rather than min-fill, and it makes explicit that the $O(1/V)$
guarantee is a total-quality guarantee. The reversal is the $V$-space
signature of the quality-fairness conflict.

\begin{figure}[!t]
\centering
\includegraphics[width=\columnwidth]{Vsweep_conflict.pdf}
\caption{The Lyapunov weight $V$ in the heterogeneous multi-satellite setting:
total quality rises with $V$ (the total-quality $O(1/V)$ improvement of
Theorem~\ref{thm_gap} holds), while min-fill falls, reversing its
conflict-free single-satellite behaviour. The reversal is the $V$-space
signature of the quality-fairness conflict.}
\label{fig:vconf}
\end{figure}

%%EVGAP_BLOCK%%

\subsection{Sensitivity Analysis}\label{EV_SENS}
We sweep the scenario parameters one at a time around the nominal point, each
value over 10 seeds (Fig.~\ref{fig:sens}).

The platform-baseline sweep exposes the 7B feasibility boundary. The peak
constraint makes the platform load a hard gate for the top tier: the 7B split
is feasible only while $P_{\text{base}} + 6.37 \le 15$\,W, that is, below about
8.6\,W. The sweep confirms the boundary empirically: at $P_{\text{base}}{=}6$\,W
the scheduler runs 842 splits with zero downtime, at 7\,W about 200, at 8\,W
the count collapses to 87, and at 9\,W the 7B tier drops out entirely; beyond
9\,W the platform load alone exceeds what the panel sustains and every policy
degrades. This boundary is set by the hardware cap and the measured 7B peak,
not by a tuned parameter, and it is the empirical basis for the per-slot peak
shield.

The battery sweep bounds the scope of the safety separation. With the battery
sized close to the eclipse requirement ($15{,}960$\,J at the nominal load),
only the proposed scheduler and the two non-splitting baselines stay within a
5\% downtime SLA while the greedy splitters black out; this deployability
partition is invariant for $B^{\max}\!\in\![16{,}000, 18{,}000]$\,J. It does
not persist when the battery is provisioned generously: at
$B^{\max}\!\ge\!19{,}000$\,J the reckless splitters no longer black out, and at
$24{,}000$\,J Random reaches a higher min-fill than the proposed scheduler
(0.367 versus 0.296) with negligible downtime. We report this boundary
explicitly; Section~\ref{EV_LIM} analyzes its cause.

The remaining axes are consistent with these two. The solar-panel sweep shows
the same energy story from the supply side (the sustainability floor sits near
$P^{\text{solar}}\!\approx\!12$\,W for the 7\,W platform); under overload
(arrival probability up to 0.7) every policy's min-fill falls while the
proposed scheduler's downtime stays at 0.4\% throughout. On the
constellation-size axis the proposed scheduler's min-fill is non-monotone at 16
satellites: an idle decomposition locates the direct cause in the
eclipse-reserve guard of Remark~\ref{rem_reserve}, whose blocked slots grow
from 423 to 1134 (of about 7960 decision slots) while score-based idling
slightly decreases; the interaction between the per-satellite reserve
requirement and the constellation phase structure is left to the
concurrent-scheduling extension.

\begin{figure*}[!t]
\centering
\includegraphics[width=\textwidth]{scope_sensitivity.pdf}
\caption{One-at-a-time sweeps around the nominal point (10-seed means, $\pm1$
std). Top row: downtime against the 5\% SLA (dashed) — who is deployable
where. Bottom row: min-fill. The battery and constellation panels show the
boundary of the safety separation honestly: with generous resources the
reckless splitters no longer black out and Random reaches higher min-fill
(Section~\ref{EV_LIM}).}
\label{fig:sens}
\end{figure*}

\subsection{Dynamic ISL Topology}\label{EV_ISL}
To make the handoff constraint~\eqref{C_pipe_isl} bind, we replace the static
topology with time-varying links of controllable availability $p_{\text{ISL}}$
(deterministic per-window connectivity, coherence window 300 slots,
Fig.~\ref{fig:isl}). As availability falls from 1.0 to 0.3 the number of 7B
splits falls from 201 to 136, so the constraint genuinely binds. The loss
concentrates entirely on the 7B-dependent source: mountain's fill drops from
0.165 to 0.154 while every other source stays flat or rises slightly as
standalone capacity flows back; min-fill falls accordingly from 0.169 to 0.157,
total quality is unchanged, and downtime stays at 0.4\%. Under heterogeneous
sources, ISL connectivity thus couples directly to the fairness of the
top-tier-dependent source, rather than to throughput or safety.

\begin{figure}[!t]
\centering
\includegraphics[width=\columnwidth]{isl_locked.pdf}
\caption{Time-varying ISL availability. (a) Feasible 7B splits fall as links
degrade, so the handoff constraint \eqref{C_pipe_isl} binds. (b) The loss lands
on min-fill through the 7B-dependent source, while total quality is unchanged.}
\label{fig:isl}
\end{figure}

\subsection{Stability Stress Tests}\label{EV_STAB}
A twenty-orbit run (114{,}000 slots) probes the guarantees directly
(Fig.~\ref{fig:stab20}). The battery settles into a periodic orbit-scale cycle,
and the energy and quality virtual queues are mean-rate stable: their growth is
sublinear and flattening, with $Q^E(T)/T \to 7.5\times10^{-4}$ and
$Q^U(T)/T \to 1.7\times10^{-3}$ (mean backlogs reaching only 86.5 and 192.3
after twenty orbits), while blackout stays below 0.1\% of satellite-slots. This
is the empirical signature of the robust stability of Theorem~\ref{thm_robust}
and the mean-rate guarantees of Theorem~\ref{thm_stable} under periodic
zero-harvest eclipse intervals.

\begin{figure}[!t]
\centering
\includegraphics[width=\columnwidth]{stability_20orbit.pdf}
\caption{Twenty-orbit stability run (proposed, $V{=}10$). (a) The battery
cycles periodically through eclipse. (b), (c) The energy and quality-deficit
virtual queues grow sublinearly and are mean-rate stable
($Q(T)/T \to 10^{-4}$--$10^{-3}$), the empirical signature of
Theorems~\ref{thm_stable} and~\ref{thm_robust}.}
\label{fig:stab20}
\end{figure}

\subsection{Limitations and Future Work}\label{EV_LIM}
Our empirical advantage is confined to the energy-constrained regime, and we
state its scope precisely, in two dimensions. First, battery capacity: the
deployability separation holds for batteries sized close to the eclipse
requirement ($B^{\max}\!\in\![16{,}000,18{,}000]$\,J) and disappears when the
battery is provisioned generously, where a policy that splits at random reaches
higher min-fill and total quality than our scheduler. An idle decomposition
verifies the cause: in the generous regime the eclipse-reserve guard of
Remark~\ref{rem_reserve} is essentially inactive (zero blocked slots at
$24{,}000$\,J), and the conservatism comes entirely from the energy virtual
queue, which tracks the energy rate deficit — the standard virtual-queue
construction that guarantees mean-rate sustainability and is correct in the
constrained regime. Because it tracks a rate deficit rather than the battery
state of charge, it does not distinguish a full battery from a tight one and
therefore throttles inference that a full battery could absorb. Folding the
state of charge into the drift, a battery-aware back-pressure term, is a
natural extension that trades a more involved frame-based analysis for
cross-regime adaptivity; we leave it to future work. Second, constellation
scale: at 16 satellites the same guard blocks about 2.7 times more decision
slots than at 8, producing a non-monotone min-fill; a concurrent
multi-satellite scheduler (the simulation commits one decision per slot, a
conservative special case of the formulation) is the corresponding extension.
Two further directions follow from the same analysis: a per-resource max-min
penalty that optimizes fairness per unit of energy rather than per task, and a
higher-fidelity power model that separates the platform power system from the
compute payload rather than treating the peak cap as one shared budget.
